\section*{Chapitre \ref{ch:b2:weather-systems} --- Systèmes météorologiques}

\begin{solution}{exo:b2:weather-systems:1}
(a) $f = 2 \times 7.292 \times 10^{-5} \times \sin\ang{50} = \qty{1.117e-4}{s^{-1}}$. (b) Avec
$\bar T = \qty{278}{K}$, $\tan\alpha = 1.117 \times 10^{-4} \times 278 \times 15/(9.81 \times 6)
= 0.00791$, soit environ 1 sur 126. (c) $1/0.00791 = 126$~: à environ \qty{126}{km} derrière le front. (d) La
pente est inversement proportionnelle au contraste de température~: elle est divisée par deux, soit 1 sur 253.
\end{solution}

\begin{solution}{exo:b2:weather-systems:2}
(a) $6 \times 150 = \qty{900}{km}$ en avant. (b) $900/25 = \qty{36}{h}$, soit un jour et demi. (c) D'abord
les cirrus élevés, puis les cirrostratus, puis les altostratus ; le nuage s'épaissit et s'abaisse, et la pluie
tombe depuis le nuage bas et épais à mesure que le front approche.
\end{solution}

\begin{solution}{exo:b2:weather-systems:3}
Les symboles indiquent le sens de déplacement du front et pointent vers l'air plus chaud devant un front froid, ou vers l'air
plus froid devant un front chaud. (a) Front froid~: l'air froid se trouve au nord-ouest, derrière lui ; il se déplace vers le sud-est.
(b) Front chaud~: l'air froid se trouve au nord, devant lui ; il se déplace vers le nord.
(c) \omterm{def:b2:weather-systems:fronts}{Front occlus}, reliant la dépression au point triple où commencent les fronts froid et chaud ; l'air froid se trouve des deux côtés,
l'air chaud étant soulevé au-dessus ; il se déplace dans le sens vers lequel pointent ses symboles.
(d) \omterm{def:b2:weather-systems:fronts}{Front stationnaire}~: les triangles pointent vers le sud, dans l'air chaud, donc l'air froid se trouve au nord ; il
bouge à peine.
\end{solution}

\begin{solution}{exo:b2:weather-systems:4}
(a) $\beta = 2 \times 7.292 \times 10^{-5} \times \cos\ang{45}/(6.371 \times 10^6) =
\qty{1.62e-11}{m^{-1}.s^{-1}}$. (b) Pour \qty{4000}{km}, $k = \qty{1.571e-6}{m^{-1}}$, $\beta/k^2 =
\qty{6.6}{m/s}$ et $c = 12 - 6.6 = \qty{5.4}{m/s}$~; pour \qty{8000}{km}, $\beta/k^2 =
\qty{26.3}{m/s}$ et $c = \qty{-14.3}{m/s}$. (c) L'onde courte se déplace vers l'est, de \qty{470}{km} par jour~;
l'onde longue se déplace vers l'ouest, de \qty{1240}{km} par jour, à contre-vent.
\end{solution}

\begin{solution}{exo:b2:weather-systems:5}
(a) $c = 0$ lorsque $U = \beta/k^2$, donc $k = \sqrt{\beta/U}$ et $\lambda = 2\pi/k =
2\pi\sqrt{U/\beta}$. (b) $2\pi\sqrt{12/(1.62 \times 10^{-11})} = \qty{5410}{km}$. (c) Le cercle
de latitude \ang{45} mesure $2\pi \times 6371 \times \cos\ang{45} = \qty{28300}{km}$~: environ 5.2
ondes, donc cinq en pratique. (d) $\lambda$ croît comme $\sqrt U$~: un vent d'ouest plus fort allonge les
ondes stationnaires, et il en tient moins autour de l'hémisphère.
\end{solution}

\begin{solution}{exo:b2:weather-systems:6}
(a) Au point $C_+ = (\sqrt{72}, \sqrt{72}, 27)$~: $\dot x = 10(\sqrt{72} - \sqrt{72}) = 0$, $\dot y =
\sqrt{72}(28 - 27) - \sqrt{72} = 0$, $\dot z = 72 - \tfrac83 \times 27 = 72 - 72 = 0$.
(b) $\mathrm{e}^{-13.67} = 1.16 \times 10^{-6}$~: le volume est presque divisé par un million en une unité de temps.
(c) Un volume peut diminuer tout en s'étirant dans une direction et en se comprimant beaucoup plus fortement dans les
autres~: la distance dans la direction d'étirement croît, et deux trajectoires voisines s'y séparent, tandis que le volume
qui les sépare s'annule.
\end{solution}

\begin{solution}{exo:b2:weather-systems:7}
(a) $\ln 2/0.906 = 0.765$ unité de temps. (b) $\ln(10/10^{-8})/0.906 = \ln(10^9)/0.906 = 22.9$ unités
de temps. (c) $\ln 1000/0.906 = 7.6$ unités de temps~: une erreur initiale mille fois plus petite ne
fait gagner qu'un tiers du temps trouvé en (b), puisque $\ln 10^3/\ln 10^9 = 1/3$.
\end{solution}

\begin{solution}{exo:b2:weather-systems:8}
(a) $(303 - 195)/195 = 0.554$ et $V_{\max} = \sqrt{0.9 \times 0.554 \times 15\,000} =
\qty{86.5}{m/s}$. (b) Bien au-dessus des \qty{70}{m/s} de la catégorie 5, proche de la valeur de \qty{86}{m/s} que
cinq tempêtes récentes ont dépassée. (c) Les tempêtes restent rarement longtemps dans l'état stationnaire que suppose la
théorie~: le cisaillement vertical du vent, l'air sec aspiré dans la tempête et la subsidence de l'air
environnant, qui ne figurent pas dans la théorie, affaiblissent les tempêtes.
\end{solution}

\begin{solution}{exo:b2:weather-systems:9}
(a) $(303 - 195)/303 = 0.356$. (b) $(303 - 195)/195 = 0.554$~; le rapport vaut $0.554/0.356 =
303/195 = 1.554$. (c) Dans un moteur de Carnot, le travail sort du moteur. Dans un cyclone tropical, le travail est
dissipé par frottement près de la surface de la mer, et cette chaleur réchauffe l'air entrant à la
température $T_s$, où le moteur prélève sa chaleur~: elle est réutilisée, et le moteur
se comporte comme s'il était plus efficace, du facteur $T_s/T_o$.
\end{solution}

\begin{solution}{exo:b2:weather-systems:10}
(a) La dérivation de la \cref{prop:b2:weather-systems:slope} sans approximation donne
$\mathrm{d}z/\mathrm{d}x = f(\rho_c v_c - \rho_w v_w)/[g(\rho_c - \rho_w)]$. Un fluide plus dense
reposant sur un fluide plus léger est instable~: la moindre perturbation le fait basculer. Dans un front au
repos, l'air froid, plus dense, doit donc se trouver dessous, comme dans le coin de la définition ; alors
$\rho_c > \rho_w$, la surface monte vers l'air froid, et $\rho_c v_c > \rho_w v_w$.
(b) $\zeta = \partial v/\partial x = 10/(50 \times 10^3) = \qty{2e-4}{s^{-1}}$, soit environ $1.9f$ avec
$f = \qty{1.03e-4}{s^{-1}}$. (c) Positive dans l'hémisphère Nord~: cyclonique, et aussi grande que
le paramètre de Coriolis lui-même ; les fronts sont des bandes de forte vorticité cyclonique.
\end{solution}

\begin{solution}{exo:b2:weather-systems:11}
(a) $\log_2(6/0.5) = \log_2 12 = 3.58$ doublements, soit $3.58 \times 2.4 = 8.6$~jours. (b) Gagner un jour de prévision
exige un doublement de moins avant que l'erreur ne devienne inutilisable~: l'erreur initiale doit donc diminuer de
$2^{1/2.4} = 1.33$ à chaque décennie. (c) Une réduction par dix correspond à
$\log_2 10 = 3.32$ doublements~; si elle ne gagne que 4 à 5 jours, les petites erreurs doublent en $4/3.32 = 1.2$
à $5/3.32 = 1.5$ jours~: les petites erreurs croissent plus vite que les erreurs plus grandes mesurées dans les
prévisions des années 1980, par la chaleur latente libérée dans la convection aux petites échelles.
\end{solution}

\begin{solution}{exo:b2:weather-systems:12}
(a) À \qty{30}{\celsius}, $e_s = \qty{43.5}{hPa}$ et $q_s^* = \qty{27.1}{g/kg}$~; l'air à
\qty{29}{\celsius} a $e = 0.8 \times 41.0 = \qty{32.8}{hPa}$ et $q_a = \qty{20.4}{g/kg}$~;
$h_o^* - h^* = 1004.5 \times 1 + 2.501 \times 10^6 \times 0.0067 = \qty{17.8}{kJ/kg}$, contre
\qty{16.8}{kJ/kg}. (b) $(303.15 - 200)/200 = 0.516$ et $V_{\max} = \sqrt{0.9 \times 0.516
\times 17\,800} = \qty{90.9}{m/s}$, contre \qty{87.9}{m/s}. (c) Environ \qty{3.4}{\percent} par
kelvin. Le déséquilibre augmente de \qty{6}{\percent}, ce qui ajoute \qty{3}{\percent} à la vitesse~;
le facteur $(T_s - T_o)/T_o$ augmente de \qty{1}{\percent}, ce qui ajoute \qty{0.5}{\percent}~: le
carburant, c'est-à-dire l'humidité que la mer plus chaude peut fournir, contribue le plus, par la croissance exponentielle de
$e_s$.
\end{solution}

\begin{solution}{pb:b2:weather-systems:1}
\textbf{1.} $(480 - 390)/6 = \qty{15}{km/h}$.
\textbf{2.} $390/15 = \qty{26}{h}$ après 12 UTC~: à 14 UTC le lendemain.
\textbf{3.} $f = \qty{1.117e-4}{s^{-1}}$.
\textbf{4.} Avec $\bar T = \qty{281}{K}$, $\tan\alpha = 1.117 \times 10^{-4} \times 281 \times
16/(9.81 \times 8) = 0.0064$, soit environ 1 sur 156.
\textbf{5.} Le front se trouve alors à $6 \times 15 = \qty{90}{km}$ au-delà de la ville ; la \omterm{def:b2:weather-systems:fronts}{surface frontale} est
à $90\,000 \times 0.0064 = \qty{576}{m}$ au-dessus d'elle.
\textbf{6.} $1500/0.0064 = \qty{234}{km}$ derrière le front, atteint $234/15 = \qty{15.6}{h}$ après
le passage.
\textbf{7.} Cumulus et cumulonimbus, averses, peut-être orages, rafales et chute brutale de
température ; la pression, en baisse à l'approche du front, se remet à monter à son passage ; puis
un air plus frais et plus clair.
\textbf{8.} $\beta = 2 \times 7.292 \times 10^{-5} \times \cos\ang{50}/(6.371 \times 10^6) =
\qty{1.47e-11}{m^{-1}.s^{-1}}$.
\textbf{9.} $k = 2\pi/(5 \times 10^6) = \qty{1.257e-6}{m^{-1}}$ et $\beta/k^2 = \qty{9.3}{m/s}$.
\textbf{10.} $c = 20 - 9.3 = \qty{10.7}{m/s}$, vers l'est.
\textbf{11.} $10.7 \times 86\,400 = \qty{924}{km}$ vers l'est.
\textbf{12.} $\lambda_s = 2\pi\sqrt{20/(1.47 \times 10^{-11})} = \qty{7330}{km}$~; le cercle de
latitude \ang{50} mesure $2\pi \times 6371 \times \cos\ang{50} = \qty{25700}{km}$, soit environ 3.5
de telles ondes.
\textbf{13.} En avant de lui, sur son flanc est, là où les ondes courtes induisent une ascendance et où la
pression de surface baisse.
\textbf{14.} $\lambda = \ln(3.9 \times 10^{-3}/5 \times 10^{-7})/(14 - 4) = \ln 7800/10 = 0.896$.
\textbf{15.} $\ln 2/0.896 = 0.77$ unité de temps.
\textbf{16.} \qty{1}{\percent} en dessous de 0.906~: la moyenne de 200 couples, un temps fini et des lectures sur
un graphique.
\textbf{17.} Une fois les deux trajectoires éloignées, elles se trouvent en des points sans lien sur l'attracteur, et
leur distance ne peut dépasser sa taille~: la croissance sature.
\textbf{18.} $\ln 10/0.896 = 2.6$ unités de temps plus tard.
\textbf{19.} $\log_2(3/0.1) = \log_2 30 = 4.91$ doublements.
\textbf{20.} $4.91 \times 2.4 = 11.8$~jours.
\textbf{21.} Un temps de doublement~: 2.4~jours.
\textbf{22.} $4.91 \times 5 = 24.5$~jours~: l'estimation initiale était beaucoup trop optimiste.
\textbf{23.} Avec un temps de doublement de 2.4 jours, une erreur dix fois plus petite gagnerait $3.32 \times 2.4
= 8.0$~jours, et non 4 à 5~: les petites erreurs croissent plus vite que 2.4 jours par doublement, et une
estimation fondée sur ce temps de doublement surestime l'horizon.
\textbf{24.} Cette prévision peut être utile au plus environ \textbf{12 jours}, et en pratique
pendant moins longtemps.
\end{solution}
